% Français 9115, batch 24 — Gallica views 461–480.
% Pass: Opus 5.5 (claude-opus-5-5), 2026-09-22 — first pass, unchecked against the views by a human.
%
% What the twenty views hold. One piece, begun before the batch and going
% on after it, in French, in her regular fair hand but more corrected
% than in batch 22 (struck words, words above the line, one long addition
% in a finer pen at v. 480): the same memoir on vibrating elastic plates
% as in batches 21–22, now in its comparison of theory with Chladni's
% experiments on square plates (§ 4, opened at v. 439), paginated by her
% (55) to (72), with paragraphs N.o 13, 14 and 15. The leaves are small,
% pasted on guards, and written on both sides: the odd pages carry the
% ink number at the top right, the even ones none, as in batch 22. One
% interloper, a torn fragment on the measure of curvature (v. 465), is
% bound between (58) and (59) and numbered in the ink series. No
% microfilm target, no printed matter, no repeated exposure.
%
%   v. 461  (55) end of the discussion of integral (C): with an odd number
%           of points of rest the lines through the middles of the sides
%           satisfy the limit conditions (« V. N.o 8 »), so the square
%           plate falls into four equal parts that may show figures of
%           different ratios scriptA : B; Chladni touching several points
%           at once forces the plate into such figures.
%   v. 462  (56) Chladni's fig. 78 (fourth value of ω): squares a n m y,
%           g k h d keep the nodal lines of her figure, the others show
%           sinuous lines; a long addition, begun above the line and
%           finished in the left margin, on the break of continuity of the
%           nodal curve; each sinuous line has two « flexions + ½ ».
%   v. 463  (57) the flexion counted (after Chladni's definition, two
%           sinuosities); fig. 85: the lower half with sinuous lines of 3
%           complete flexions through six invariable points of rest; then
%           the experiment she made herself: on a square plate showing
%           Chladni's fig. 64 (two diagonal nodal lines) —
%   v. 464  (58) — she had the quarter B C D E cut away, fingers kept at c,
%           bow at F: same sound, same figure (« Voyez ci-joint fig. A »,
%           not on the leaf); so ends free by dz/dx = 0, d³z/dx³ = 0, which
%           she could not verify on the simple lamina (« V. § 3 N.o 8 »),
%           are confirmed.
%   v. 465  a fragment, not part of the memoir: torn at the top, larger
%           hand, no pagination, ink number 228. « je proposerai d'abord a
%           l'examen des g[…] les deux théorèmes suivants »: the curvature
%           of a surface measured by the sum of the inverses of the
%           principal radii; its irregularity by their difference.
%   v. 466  the blank back of that fragment: not transcribed.
%   v. 467  (59) the memoir resumes mid-sentence from (58): why only
%           rectangular surfaces; diagonal nodal lines satisfy z = 0 and
%           d²z/dx² + d²z/dy² = 0; right isosceles triangles with sides
%           supported and hypotenuse free « au sens d'Euler » —
%   v. 468  (60) — and the converse cases, all contained in the equations
%           of the square plate, changing ω into ½ω; the section closes
%           mid-page.
%   v. 469  (61) homage to « l'étonnante sagacité de M^r Chladni », who
%           judged the bizarre figures equivalent to parallel straight
%           lines; N.o 13 « Surfaces périodiques »: the integrals of N.o 4
%           and 5, z = … sin(π nx/A) sin(π my/A), the plate supported all
%           round.
%   v. 470  (62) « ligne d'appui » defined; a period between four lines of
%           support, cut into quarters by analytic limits; the square
%           supported on four sides and the square free on four sides
%           governed by the same equation —
%   v. 471  (63) — the quarters reversed; the likeness with the string and
%           the open flute (« V N.o 7 »); three lines hatched out; n and m
%           straight nodal lines, equally spaced, halved by analytic
%           limits.
%   v. 472  (64) every intersection of the lines ending the periods may be
%           an origin (N.o 9); integral (D) of N.o 4 admits no other nodal
%           lines; Chladni's figs 63, 66a, 68a, 71a, 74a, 75a, 80a, 81a and
%           82 (« V. T. d'acoustique ») belong to it; fig. 63, m = n = 1.
%   v. 473  (65) figs 63, 66a, 74a, sound proportional to m² + n² = 2, 5,
%           17 (« V. N.o 4 »); Chladni's « vibrations tournantes » (« V. T.
%           d'acoustique N.o 84 … 87 et N.o 98 ») —
%   v. 474  (66) — differ only in having one period in x; figs 68a, 71a,
%           75a, m² + n² = 8, 13, 18.
%   v. 475  (67) figs 80a, 81a, 82, m² + n² = 25, 29, 32; N.o 14: the
%           series factor of integrals (E) and (F) of N.o 4, constants a′,
%           b′, A′, B′ free, one sound for several figures —
%   v. 476  (68) — sounds (Π−1)(m² + n²) for (E), m²(Π+Π′−2) for (F);
%           with Π = 1 − p², Π′ = 1 − q², of the same form as for (D).
%   v. 477  (69) with b′ = B′ = 0 or a′ = A′ = 0 the series of (E) becomes
%           sin(π npx/A), cos(π npx/A), « formule P. 143 des leçons sur le
%           calcul des fonctions in.8 an 1806 »; the same for the second
%           family, written « (E) » on the leaf.
%   v. 478  (70) no other case terminates the series; she leaves them;
%           N.o 15: integral (G) of N.o 5, reducing to (F), or to z″ the sum
%           of two separate functions of x and y —
%   v. 479  (71) — the free square plate needs a″ = A″ = 0; with Π″ = 1−9,
%           z = … sin sin {b cos² + B cos² − (b+B)/4}; Chladni's figs 69,
%           70, 87a, 87b, 88a, 88b; fig. 69 (b = −B, two diagonals), fig.
%           70 (a relation between b and B she could not measure).
%   v. 480  (72) four nodal lines placed by cos²(πx/A) = (b+B)/4b; the
%           case b = −B; a long addition in a finer pen on the cases where
%           a condition is impossible, where b or B vanishes (straight
%           lines through cos(πx/A) = ½, the « lignes fictives » of (C),
%           « véritables points invariables de repos »); « Dans tout autre
%           cas » four inclined nodal lines. The page ends on a full stop
%           two thirds down; v. 481 (ink 236, page (73)) opens « On voit
%           donc que la théorie est parfaitement d'accord avec
%           l'expérience … fig. 70 ».
%
% Seams. View 461 opens mid-sentence (« … se / rapporteroient à la même
% valeur de l'angle ω ») on page (55): pages (35)–(54) and the ink
% numbers 216–225 fall in batch 23 (v. 441–460, not read here; view 460
% is not in the mirror). View 480 ends a paragraph on a full stop; view
% 481 continues the same discussion of fig. 70 on page (73). References to
% paragraphs outside the batch: N.o 4, 5 (the integrals, v. 405–409,
% batch 21), N.o 7 (v. 414), § 3 N.o 8, N.o 9 (v. 433), and « l'intégrale
% (C) » of N.o 3.
%
% The numbers on the leaves, and why no \folio{} is written. (1) Her
% pagination, at the head, centre: (55) at v. 461, (56)–(64) at v.
% 462–472 without the fragment, (65)–(72) at v. 473–480, without a gap;
% the odd pages (55), (57), … (69) are struck with an oblique stroke, the
% even pages not; page 71 (v. 479) is written without parentheses and
% struck. (2) The ink series at the top right, in the larger rounder
% hand: 226 (v. 461), 227 (v. 463), 228 (v. 465, the fragment), 229 (v.
% 467), 230 (v. 469), 231 (v. 471), 232 (v. 473), 233 (v. 475), 234 (v.
% 477), 235 (v. 479) — one per physical leaf, without a gap, on the side
% paginated odd; 236 stands on v. 481. It counts the fragment as a leaf,
% as batch 22 saw it count a slip (210, v. 428). Penned, not pencilled:
% no \folio{}. (3) Her paragraph numbers: N.o 13 (v. 469), N.o 14 (v.
% 475), N.o 15 (v. 478). (4) Chladni's figure numbers (63–88b) and his
% paragraph numbers (N.o 84–87, 98), cited. Every number above was read on
% its view; none is inferred. None of these leaves can be Laubenbacher and
% Pengelley's « Manuscript C » (ff. 348r–349r): the subject is the elastic
% plate, and the ink series here runs 226–235.
%
% Notation. Hers, kept. The first coefficient of (C) at v. 461 is the
% looped cursive capital of v. 439–440, set $\mathcal{A}$; A′, B′, A″, B″
% of the later integrals are upright capitals. The tall-legged capital of
% the sound formulas and of Π = 1 − p² (v. 476–479) is set Π, Π′, Π″ as
% in batch 21; the π of sin(π nx/A) is set π. Her « sin », « cos » are
% written with a long s and set \text{sin}, \text{cos}; the points of
% suspension « z = …… » standing for an omitted factor are hers. Words
% and letters she underlines in the prose are set in \emph{}; letters
% underlined inside mathematics keep \underline{}. Her 4 is written like
% « h » (« 64 », « 74a », « N.o 4 », « N.o 84 »). Spelling is hers and is
% not flagged: « quarré(e) », « parallelle », « differens », « mouvemens »,
% « appartenans », « indépendans », « coëfficiens », « complette »,
% « emploirai », « emploiées », « fesant », « tracés » for tracées; the
% run-together words « dela », « ala », « àla »,
% « àlafois », « aucontraire », « desorte », « demême », « parconséquent »,
% « cequi », « ceque », « aunombre », « devibrer », « adire ». Where the
% leaf and the calculation or the sense disagree the leaf stands, with a
% note: v. 467 (a derivative written without its denominator dx), v. 470
% (dz/dy without « = 0 », a condition ending on « = »), v. 477 (« (E) »
% where the constants are those of (F)), v. 479 (« Π″ = 1 − 9 »).
%
% What could not be read, and what is offered with doubt. \ill{} stands
% six times: the torn first word of the fragment (v. 465), a struck word
% after « pour » (v. 468), the label of an integral in three hatched lines
% (v. 471), the end of a struck line (v. 473), a word in a struck line (v.
% 479), a word in a struck line after « l'origine » (v. 480). \uncertain{} stands for « tel », « naturel », « en eff » (v.
% 461); « dans », « qui », « -sent », r (v. 462); fig. « 85 », « lignes »,
% « tasse », fig. « 64 » (v. 463); « on sait », « m. », fig. « A » (v.
% 464); « -ts » (v. 465); « avantage » (v. 468); « homage » (v. 469);
% « présenté », « l'observation », « s'applique », « satisfera » (v. 471);
% « cette » (v. 473); « duquel », « telle » (v. 475); « vient »,
% « celui » (v. 476); « fait », « soit », « éd. » (v. 477); « qu », « p »
% (v. 479); a struck group « B par exemple sera > que 3b » (v. 480).
% Nothing was guessed to fill a gap.
%
% Nothing here was checked against any published edition of Chladni, of
% Lagrange's Leçons, or of Sophie Germain's papers.
\documentclass[11pt,a4paper]{article}
% The preamble shared by every transcription of the mathematicians' manuscripts in Gallica.
%
% It rests on one idea: **uncertainty compounds, it does not resolve.** A
% transcription that smooths over a doubtful word has destroyed the only thing
% it was made for — a reader must be able to tell, at every point, what was on
% the page and what was supplied. The apparatus macros below are therefore the
% critical apparatus, not decoration.
%
% The same commands are recognised by `scripts/render.mjs`, which produces the
% site's reading view, and by `scripts/tei.mjs`. A macro added here must be
% added there too, or rendering fails loudly — which is the wanted behaviour.
%
% Comments here are English, like the rest of the repository. The documents
% this preamble sets are French, and that is deliberate: see the skills.

\ProvidesPackage{germain}[2026/09/15 Transcriptions of mathematicians' manuscripts in Gallica]

\usepackage{amsmath,amssymb,amsthm}
\usepackage{mathrsfs}
% Kept from the parent preamble so the renderer's diagram support stays
% available; these manuscripts draw no commutative diagrams, and nothing here uses it.
\usepackage{tikz-cd}
\usepackage[english,french]{babel}
\usepackage{fontspec}

% --- Greek ------------------------------------------------------------------
% The text face stays fontspec's default, Latin Modern, which has no Greek:
% XeTeX drops every glyph it lacks with a line in the log and nothing on the
% page, and the Greek words the manuscripts quote vanished from the PDFs. So
% the Greek and Greek Extended blocks get a XeTeX character class of their own,
% and entering or leaving a run of it switches the family to CMU Serif — the
% same Computer Modern design, with polytonic Greek — and back. Only the
% family changes: a Greek word in \textit or \textbf keeps its shape and series.
% The transitions to and from every other class carry the tokens that class
% has towards an ordinary letter, so French spacing before « ; » or inside
% « » still holds after a Greek word. No grouping, so no brace or \endgroup
% can come out unbalanced; maths is left alone, since XeTeX inserts nothing
% there. Kept identical in `exercices.sty`.
\makeatletter
\newfontfamily\germainGreekfont{cmun}[
  Extension=.otf, NFSSFamily=germaingreek,
  UprightFont=*rm, ItalicFont=*ti, SlantedFont=*sl,
  BoldFont=*bx, BoldItalicFont=*bi, BoldSlantedFont=*bl]
\newXeTeXintercharclass\germain@greek
\def\germain@greekrange#1#2{%
  \@tempcnta=#1\relax
  \loop
    \XeTeXcharclass\@tempcnta=\germain@greek
    \ifnum\@tempcnta<#2\advance\@tempcnta\@ne\repeat}
\germain@greekrange{"0370}{"03FF}
\germain@greekrange{"1F00}{"1FFF}
\def\germain@greekprev{\rmdefault}
\def\germain@greekin{%
  \edef\germain@greekprev{\f@family}\fontfamily{germaingreek}\selectfont}
\def\germain@greekout{\fontfamily{\germain@greekprev}\selectfont}
\def\germain@greekjoin{%
  \XeTeXinterchartokenstate=\@ne
  \@tempcnta=\z@
  \loop
    \ifnum\@tempcnta=\germain@greek\else
      \XeTeXinterchartoks\germain@greek\@tempcnta=\expandafter{%
        \expandafter\germain@greekout\the\XeTeXinterchartoks\z@\@tempcnta}%
      \XeTeXinterchartoks\@tempcnta\germain@greek=\expandafter{%
        \the\XeTeXinterchartoks\@tempcnta\z@\germain@greekin}%
    \fi
    \ifnum\@tempcnta<4095\advance\@tempcnta\@ne\repeat}
% babel-french sets its own classes, and saves and restores the interchar
% state, each time a language is selected: join again after it does.
\addto\extrasfrench{\germain@greekjoin}
\addto\extrasenglish{\germain@greekjoin}
\AtBeginDocument{\germain@greekjoin}
\makeatother

\usepackage{geometry}
\geometry{a4paper,margin=25mm,marginparwidth=28mm}
\usepackage{marginnote}
\usepackage{xcolor}
\usepackage[normalem]{ulem}
\setlength{\emergencystretch}{3em}
\usepackage{hyperref}
\hypersetup{colorlinks=true,linkcolor=black,urlcolor=black,pdfborder={0 0 0}}

% --- Batch metadata ---------------------------------------------------------
% Taken as they stand by the rendering script, which shows them at the head of
% the reading view. They fix what the file claims to cover. The macro names are
% the parent project's, so that the scripts stay shared; \folder here means the
% volume's slug — `fr-9115` — and \pages counts Gallica views.

\newcommand{\folder}[1]{\gdef\grothFolder{#1}}
\newcommand{\batch}[1]{\gdef\grothBatch{#1}}
\newcommand{\pages}[2]{\gdef\grothFirst{#1}\gdef\grothLast{#2}}
\newcommand{\foldertitle}[1]{\gdef\grothFoldertitle{#1}}

% The holder's own shelfmark, printed as they write it: « Français 9115 ».
\newcommand{\shelfmark}[1]{\gdef\grothShelfmark{#1}}

% The dating, copied verbatim from the holder's catalogue — never inferred
% here. « XIXe siècle » is the BnF's; a pass that dates a leaf from its content
% says so in a \note{}, as its own inference.
\newcommand{\dating}[1]{\gdef\grothDating{#1}}

% The status notice, printed in the title block and repeated as a diagonal
% watermark on every page of the PDF. The manuscripts are in the public
% domain; what this line declares is not a right but a quality — an
% unverified first pass by a machine. The skills write the line; a file
% without it gets no watermark, so leaving it out is visible in the output.
\newcommand{\watermark}[1]{\gdef\grothWatermark{#1}}

\gdef\grothFolder{?}\gdef\grothBatch{}\gdef\grothFirst{?}\gdef\grothLast{?}%
\gdef\grothFoldertitle{}\gdef\grothShelfmark{}\gdef\grothDating{}\gdef\grothWatermark{}

% --- The résumé -------------------------------------------------------------
\newenvironment{resume}
  {\small\setlength{\parskip}{0.45em}}
  {\par\medskip\hrule\medskip}

% --- Keywords ---------------------------------------------------------------
% In English, closing the modernised reading's résumé: the modern vocabulary
% under which to search for what the leaves construct. The manifest extracts
% them to tag the volume — this line is the single source of the tags.
\newcommand{\keywords}[1]{\par\smallskip\noindent
  {\footnotesize\textsc{Keywords} --- \textit{#1}}\par}

% --- The critical apparatus -------------------------------------------------

% The Gallica view number, in the margin. This is the anchor that lets a
% disagreement over a formula be settled by looking at *one* image rather than
% a volume of 750. The site uses it to turn the facsimile as the transcription
% is scrolled. A view is one scanned image: usually one side of a leaf.
\newcommand{\page}[1]{%
  \marginnote{\footnotesize\color{gray!70}\textsf{v.\,#1}}%
  \label{page:#1}\ignorespaces}

% The BnF's pencilled foliation, where the leaf carries it and a pass can read
% it: « 198r ». This is what the literature cites — « ff. 198r–208v » — and
% the one line that turns a citation into a view. Recorded, never inferred:
% a folio a pass cannot read is not written.
\newcommand{\folio}[1]{%
  \marginnote{\footnotesize\color{gray!50}\textsf{f.\,#1}}\ignorespaces}

% The modernised reading's coarser counterpart to \page{}: which views a
% section is drawn from.
\newcommand{\pagerange}[2]{%
  \marginnote{\footnotesize\color{gray!70}\textsf{v.\,#1--#2}}%
  \ignorespaces}

% Illegible: never guessed.
\newcommand{\ill}{\textcolor{red!60!black}{[\,\ldots\,]}}

% A doubtful reading: the word is given, and flagged as uncertain.
\newcommand{\uncertain}[1]{\uline{#1}}

% An editorial addition — a missing letter, an implied word.
\newcommand{\add}[1]{\textcolor{blue!50!black}{[#1]}}

% Struck out by the writer. What was crossed out is part of the document.
\newcommand{\struck}[1]{\sout{#1}}

% The transcriber's note, on the state of the page or the mathematics.
\newcommand{\note}[1]{\par\smallskip{\small\color{gray!60!black}\textit{[#1]}}\par\smallskip}

% A marginal note of hers, distinct from the body of the text.
\newcommand{\marginal}[1]{\par\smallskip\noindent\hspace{1em}%
  {\small\color{gray!60!black}#1}\par\smallskip}

% --- Title ------------------------------------------------------------------
% The title block is French, like the documents themselves.

\AddToHook{shipout/background}{%
  \ifx\grothWatermark\empty\else
    \begin{tikzpicture}[remember picture, overlay]
      \node[rotate=52, scale=3.4, align=center, text=gray!18,
            font=\sffamily\bfseries]
        at (current page.center) {\grothWatermark};
    \end{tikzpicture}%
  \fi}

\AtBeginDocument{%
  \begin{center}
    {\large\bfseries Manuscrits de mathématiciens (Gallica) ---
      \ifx\grothShelfmark\empty\grothFolder\else\grothShelfmark\fi}\\[2pt]
    {\itshape\grothFoldertitle}\\[4pt]
    {\small vues \grothFirst--\grothLast
      \ifx\grothBatch\empty\else\ (lot \grothBatch)\fi}\\[2pt]
    \ifx\grothDating\empty\else
      {\small\color{gray!60!black}Datation du catalogue : \grothDating}\\[2pt]
    \fi
    \ifx\grothWatermark\empty\else
      {\small\bfseries\color{red!45!gray}\grothWatermark}\\[2pt]
    \fi
    {\footnotesize\color{gray!60!black}Transcription établie d'après les images IIIF de Gallica.
     Source gallica.bnf.fr / Bibliothèque nationale de France.\\
     Les lectures incertaines sont soulignées, les illisibles notées [\,\ldots\,],
     les ajouts entre crochets ; \textsf{v.} numérote les vues de Gallica,
     \textsf{f.} la foliotation au crayon de la BnF.}
  \end{center}
  \vspace{1em}}

\endinput


\folder{fr-9115}
\shelfmark{Français 9115}
\batch{24}
\pages{461}{480}
\dating{1801-1900}
\watermark{Lecture automatique\\non vérifiée}
\foldertitle{Papiers de Sophie GERMAIN. — Recueil de dissertations et problèmes mathématiques et physiques.}

\begin{document}

\note{Numérotation des feuillets. Chaque page porte en tête, au milieu,
sa pagination : « (55) » à la vue 461, « (56) » à « (64) » aux vues 462
à 472, « (65) » à « (72) » aux vues 473 à 480, sans lacune ; les pages
impaires, de (55) à (69), sont barrées d'un trait oblique, les paires
non ; la page 71 (vue 479) est écrite sans parenthèses et barrée. Les
feuillets sont écrits des deux côtés, et seul le côté de pagination
impaire porte en haut à droite, à l'encre, d'une écriture plus ronde, un
nombre de la série qui court dans tout le volume : 226 (vue 461), 227
(463), 229 (467), 230 (469), 231 (471), 232 (473), 233 (475), 234 (477),
235 (479) ; le fragment de la vue 465, étranger au mémoire, porte 228.
Cette série numérote des feuillets et non des pages, et elle compte le
fragment ; elle occupe la place de la foliotation de la Bibliothèque,
mais elle est tracée à la plume et non au crayon, et aucun « f. » n'est
donc porté. Tous ces nombres ont été lus sur la vue ; aucun n'est déduit.
La vue 466, dos blanc du fragment, n'est pas transcrite.}

\page{461}

rapporteroient àla même valeur de l'angle $\underline{\omega}$; mais elles \struck{pour}
pourroient présenter des figures differentes entr'elles, si on établissoit,
pour chacune d'elles, des rapports differens entre les constantes
$\mathcal{A}$ et $B$. La même chose a donc lieu ici; et parceque \struck{\uncertain{tel}}
dans le cas d'un nombre impair quelque soit la valeur de
l'angle $\underline{\omega}$, le point du milieu dela simple lame satisfait
toujours aux conditions des limites (V. N.o 8), les lignes que l'on
peut imaginer par le milieu de chacun des côtés dela
plaque quarrée, satisferont aussi aux mêmes conditions
en n'ayant égard qu'à une seule des variables $\underline{x}$ et $y$;
et cette plaque se trouvera partagée en quatre portions
égales qui pourront présenter chacune des figures
dependantes de rapports differens entre $\mathcal{A}$ et $B$.

Lorsque l'on fait attention àla manière dont M\textsuperscript{r}
Chladni a formé les figures qu'il a obtenues,
il semble tout \uncertain{naturel} \struck{en effet} que les differences
dont il s'agit, se realisent \add{\uncertain{en eff}}; \add{car} cet \struck{in} \add{habile physicien} prescrit de
toucher àlafois plusieurs des points de ces figures \add{nodales};
\struck{si} s'il arrive\struck{nt} donc que les points de \emph{touchement}
\struck{appartiennent} qui se trouvent dans chacune des portions
separés par des lignes de limites, appartiennent à des
figures differentes, la plaque sera forcée de se ployer
dans ses differentes parties aux figures \struck{auxquelles peut appartenir}
\note{En tête, au milieu, « (55) », barré ; en haut à droite, à l'encre,
d'une écriture plus ronde, « 226 ». Le feuillet est petit, collé sur un
onglet : le texte s'arrête au bas du feuillet, aux deux cinquièmes de la
vue ; au-dessous, la feuille de montage, blanche, porte dans le coin
inférieur droit deux ou trois traits obliques au crayon, qui ne sont pas un
nombre. La première ligne achève une phrase commencée hors de ce lot
(« … se / rapporteroient »). Le premier coefficient est la capitale
cursive bouclée, distincte du $A$ du côté de la plaque, rendue
$\mathcal{A}$ comme aux vues 439–440. « (V. N.o 8) » renvoie à un
paragraphe qui n'est pas dans ce lot. Après « parceque », un mot court,
lu « tel », paraît barré. La fin de « naturel » est chargée de traits
et pourrait porter un « le » barré ; « en effet » est barré après lui.
Après « realisent », quelques lettres fines, lues avec doute « en eff »,
sont écrites en petit au niveau de la ligne, suivies d'un point-virgule ;
« car » est écrit au-dessus ; après « cet », un mot commencé (« in ») est
barré et « habile physicien » écrit au-dessus. « nodales » est ajouté en
petit en fin de ligne. « si » est barré devant « s'il », et la fin
« -nt » de « arrivent » est barrée. « touchement » est souligné.
« auxquelles peut appartenir » est barré en fin de feuillet ; dans
« peut », une lettre suivante est aussi barrée. La phrase se poursuit à
la vue suivante.}

\page{462}

\struck{pour} \add{\uncertain{dans}} lesquelles les points de \emph{touchemens} appartiennent aux lignes
nodales. Dans la figure 78 qui se rapporte à la
quatrième valeur de l'angle $\underline{\omega}$, \struck{le} Ce sont les parties dela
figure comprise dans les quarrés $a\,n\,m\,y$, $g\,k\,h\,d$, \struck{\uncertain{qui}} qui \struck{lesquelles}
\struck{se} réunit\struck{\uncertain{sent}} \struck{a l'angle $y$} \struck{ou} au point $y$, qui conservent les mêmes
lignes nodales que dans la figure que j'ai \struck{de} trouvée;
appartenir a la condition $\mathcal{A} = -B$.+ \add{car la solution de continuité de la courbe nodale dans les points qui repondent}
\marginal{a ceux que j'ai marqués \uncertain{$\underline{r}$} et $\underline{x}$ ne peut être
qu'accidentelle, et on en rencontre bien des exemples
semblables dans d'autres cas.} \struck{Mais} Dans les deux
autres quarrés les courbures ponctuées sont remplacées par
deux lignes sinueuses. \struck{et} il se trouve aussi dans ces quarrés
une partie courbe $i\,i$ que j'ai ponctuée quoiqu'elle
appartienne a la fig. que j'ai décrite parcequ'elle
peut en même tems être considérée comme la partie
d'une troisième ligne sinueuse parallelle aux deux
autres qui, si la plaque entière vibroit en présentant
la même figure, continueroit son cours dans une
direction parallelle aux deux premières.

J'ai dit plus haut que, pour la quatrième valeur
de l'angle $\omega$, lorsque la figure nodale étoit
composée de lignes sinueuses, chacune de ces lignes
devoit presenter deux \emph{flexions} $+\frac{1}{2}$. On doit donc avoir
\struck{pa} dans les portions de lignes sinueuses qui appartiennent
a deux des quarts de la plaque quarrée une \emph{flexion} $+\frac{1}{4}$
\struck{l'on voit en effet}, en se rapellant qu'une \emph{flexion} est
\note{En tête, au milieu, « (56) », non barré ; aucun nombre en haut à
droite : c'est le dos du feuillet de la vue 461. Le feuillet, petit, ne
couvre que les deux tiers supérieurs de la vue ; au-dessous, la feuille
de montage. La première ligne continue la phrase de la vue 461 (« aux
figures / … lesquelles ») : « pour » est barré et un mot court, lu avec
doute « dans », est écrit au-dessus. « touchemens », « flexions »,
« flexion » sont soulignés. « figure 78 » renvoie, semble-t-il, aux
figures de Chladni ; les lettres $a\,n\,m\,y$, $g\,k\,h\,d$, $i\,i$
désignent des points d'une figure qui n'est pas sur ce feuillet. Devant
« Ce sont », un « le » est barré, et le C, très grand, est écrit par-dessus.
Au-dessus de « qui », un « qui » plus petit est barré ; « lesquelles » est
barré en fin de ligne ; à la ligne suivante, « se » est barré, la fin de
« réunit » est surchargée et barrée (lue avec doute « -sent »),
« a l'angle y » est barré, et « au » est écrit sur un « ou » barré.
L'addition « car la solution de continuité … qui repondent » est écrite
au-dessus de la ligne « appartenir a la condition », précédée d'une
croix, et elle est placée ici au point où la même croix l'appelle, après
« $-B$. » ; elle se poursuit dans la marge de gauche (« a ceux que j'ai
marqués … dans d'autres cas. »), donnée ici comme note marginale. Dans la marge, les deux lettres
soulignées sont lues avec doute $r$ et $x$. « que j'ai trouvée ; / appartenir a la
condition » est gardé tel qu'il est écrit, point-virgule compris. « Mais » est barré et le D de
« Dans » repassé ; « et » est barré devant « il se trouve ». Au début de
la dernière ligne, « l'on voit en effet » est barré. La phrase se
poursuit à la vue suivante.}

\page{463}

composée, d'après \struck{m\textsuperscript{r} Chlad} la définition de M\textsuperscript{r} Chladni
\struck{d'une distance} d'un éloignement et d'un raprochement de
la ligne droite qu'il suppose c'est adire, comme je l'ai déjà
remarqué de deux sinuosités l'une superieure et l'autre inferieure
ala ligne fictive\struck{s}, \add{on voit en effet} que les lignes dont il s'agit ici, presentent
chacune une \emph{flexion} complette formée des sinuosités entre la troisième
et la quatrième, et entre la quatrième et la cinquième ligne\struck{s} fictive\struck{s},
et d'un quart de flexion pris entre la cinquième ligne fictive
et le bord de la plaque.

Al'égard de la fig. \uncertain{85}, il est visible que la moitié superieure
est entièrement conforme à celle que j'ai décrite pour ce cas
dans la supposition $\mathcal{A} = -B$; et que \add{dans} la moitié inferieure
\struck{presente dans laquelle} les \uncertain{lignes} nodales ponctuées qui conviendroient
àla même supposition, sont remplacées par des lignes sinueuses
qui, conformément àla théorie, presentent 3 flexions complettes,
et passent parconséquent chacune par les six points de repos
invariables situés sur la ligne fictive qui lui sert d'axe.
\add{la distance entre les points $c$ et $s$ ne peut être sensible parceque le point $c$ appartenant
àla fois a deux lignes nodales, la poussière se \uncertain{tasse} autour de lui}

Quoique je regarde comme suffisamment prouvé que les
lignes \struck{qui} par lesquelles on\struck{t} peut concevoir la plaque quarrée\add{, dans le cas de l'intégrale (C),}
partagée en quatre portions égales, sont, des limites
analytiques parfaitement equivalentes à des extrêmités
effectives, j'ai desiré confirmer encore \struck{le} ce résultat, en le
soumettant àl'experience. Pour cela, après avoir formé
sur une plaque quarrée la fig \uncertain{64} de M\textsuperscript{r} Chladni, laquelle
\note{En tête, au milieu, « (57) », barré ; en haut à droite, à l'encre,
« 227 ». Le feuillet est collé bas sur sa feuille de montage : le premier
tiers de la vue est blanc. La première ligne achève la phrase de la
vue 462 (« qu'une flexion est / composée »). « m\textsuperscript{r}
Chlad » est barré avant « la définition » ; « d'une distance » est barré
en tête de la deuxième ligne. « on voit en effet » est écrit au-dessus
de « fictive », dont le « s » final est barré ; de même aux « lignes
fictives » plus bas. Devant « flexion », souligné, « une » est surchargé.
Le premier mot de l'alinéa « Al'égard » a son initiale empâtée ; le
numéro de la figure se lit 85, peut-être 89. « dans » est écrit
au-dessus de « la moitié » ; à la ligne suivante, « presente » et « dans
laquelle » sont barrés, et le mot qui suit « les » est surchargé (lu
« lignes »). L'addition « la distance entre les points c et s … autour
de lui » est écrite en petit entre la fin de l'alinéa et « Quoique »,
sans signe d'appel ; les lettres $c$ et $s$ renvoient à la figure, qui
n'est pas sur le feuillet. « dans le cas de l'intégrale (C), » est écrit
au-dessus de « la plaque quarrée ». « ce » est écrit sur « le » devant
« résultat ». Le numéro de la figure de Chladni se lit 64, le 4 tracé
comme un « h ». La phrase se poursuit à la vue suivante.}

\page{464}

appartient ala première valeur de l'angle $\underline{\omega}$, et est composée
comme \uncertain{on sait} de deux lignes diagonales nodales, j'ai fait \struck{le} couper la
partie $B\,C\,D\,E$, c'est adire le quart de la plaque, en fesant
pourtant reserver l'espace $\underline{c}$ auquel les doigts avoient été
appliqués, pour former la fig 64 sur la plaque entière.

J'ai de nouveau replacé les doigts en $\underline{c}$, et j'ai continué
\struck{a tirer le son} d'appliquer l'archet en $F$ comme je l'avois
fait sur la plaque entière. Le son est resté le même
que sur la plaque complette, et la figure est demeurée
aussi nette que sur cette même plaque. \struck{seulement de sorte
que l'on peut concevoir la figure} + \uncertain{m.} Les lignes pleines terminent
la partie de la plaque que j'ai conservée \struck{et les ligne} et la
figure nodale qui s'y est formée; tandis que les lignes ponctuées
indiquent le contour de la partie enlevée et la figure
nodale qui s'y seroit formée si elle eut continué a
faire partie dela plaque vibrante (Voyez ci-joint fig. \uncertain{A})

La première valeur de l'angle $\underline{\omega}$ est aunombre des cas où
les points de milieu sont libres en vertu des conditions
$\frac{dz}{dx}=0$ et $\frac{d^3z}{dx^3}=0$. On voit donc que l'existence des extrêmités
soumises à de pareilles conditions que je n'avois pû
constater par des experiences directes sur la simple lame
(V. § 3 N.o 8) reçoit ici une entière confirmation.

L'intégrale (C) dont je viens d'examiner les conséquences, ne peut
s'appliquer dans toute sa généralité qu'aux seules surfaces rectangulaires
\note{En tête, au milieu, « (58) », non barré ; aucun nombre en haut à
droite : c'est le dos du feuillet de la vue 463, collé bas sur sa feuille
de montage comme lui. La première ligne achève la phrase de la vue 463
(« la fig 64 de M\textsuperscript{r} Chladni, laquelle / appartient »).
Après « comme », deux mots courts et serrés sont lus avec doute « on
sait ». « le » est barré devant « couper ». « 64 », après « fig »,
est serré contre « sur » ; le 4 est tracé comme un « h », comme à la vue
463. « $c$ » est souligné. Au début de la ligne « d'appliquer
l'archet », « a tirer le son » est barré. « seulement de sorte / que l'on
peut concevoir la figure » est barré sur deux lignes ; au-dessus de la
fin, une croix et quelques lettres, lues avec doute « m. », semblent un
signe d'appel sans texte correspondant sur ce feuillet. Après
« conservée », une petite marque en forme de $\neq$ sous la ligne ;
« les » est surchargé en « la » et « ligne » est barré. Le renvoi « fig.
A » (lu aussi « fig. 1 ») vise une figure jointe qui n'est pas sur ce
feuillet. Le renvoi « (V. § 3 N.o 8) » est lu sur un signe de
paragraphe tracé comme un 8 bouclé ; il renvoie à un paragraphe qui
n'est pas dans ce lot. La phrase se poursuit à la vue suivante.}

\page{465}

\ill{}ment je proposerai d'abord a l'examen des g
les deux théorèmes suivan\uncertain{ts}.

La courbure d'une surface a pour mesure la
somme des raisons inverses des deux principaux
rayons de courbure.

L'irrégularité dela courbure d'une surface a
pour mesure la difference des raisons inverses
de ses deux rayons principaux de courbure
\note{Un fragment isolé, d'un autre sujet que les feuillets qui
l'entourent : la mesure de la courbure des surfaces. C'est un morceau de
feuillet déchiré en haut, collé sur une feuille de montage, écrit d'une
écriture plus grande et plus appuyée, sans rature ; il ne porte pas de
pagination. En haut à droite, à l'encre, d'une écriture plus ronde,
« 228 », à la suite du 227 de la vue 463 : la série numérote ce fragment
comme un feuillet. La déchirure du haut emporte une ligne entière, dont
il ne reste que le bas de quelques lettres, et le début de la première
ligne conservée, dont on lit seulement la fin d'un mot, « -ment » (le
reste, sur la lèvre de la déchirure, paraît barré). Le dernier mot de
cette ligne s'arrête sur « g », sans suite visible sur le fragment ;
rien n'est complété. « suivants » : la fin du mot est surchargée. La
dernière ligne finit sans point ; le reste du fragment est blanc. La vue
466 en montre le dos, blanc, où ne paraît que l'écriture de cette face
vue par transparence ; elle n'est pas transcrite.}

\page{467}

car ces surfaces sont en effet les seules pour les points d'extrêmités desquelles
une des variables puisse \add{toujours} être regardée comme constante, desorte que
les conditions $\frac{dz}{dx}=0$ et $d.\left\{\frac{ddz}{dx^2}+\frac{ddz}{dy^2}\right\}=0$\add{, par exemple,}, qui conviennent ala
plaque soient parfaitement equivalentes à celles-ci: $\frac{dz}{dx}=0$ et $\frac{d^3z}{dx^3}=0$
qui appartiennent a la simple lame.

Cependant il peut exister sur la plaque quarrée \struck{qui se meut}
dont les mouvemens sont compris dans l'intégrale (C), d'autres lignes
de limites que celles qui passent par le milieu de cette plaque
et sont parconséquent parallelles aux axes des $\underline{x}$ ou $\underline{y}$.

Les diagonales nodales, par exemple, jouissent de cette propriété;
car elles remplissent les conditions $z=0$ et $\frac{ddz}{dx^2}+\frac{ddz}{dy^2}=0$

\struck{Ainsi s'il} \add{Si donc il} s'agit des cas où les valeurs de l'angle $\underline{\omega}$ donnent
un nombre pair de nœuds sur la simple lame, on peut
conclure que les mouvemens d'une surface triangulaire isocèle
dont les deux côtés adjacens àl'angle droit seroient appuyés
au sens d'Euler et l'hypothénuse libre au sens du même
géomètre, seront compris, pour les mêmes valeurs de $\underline{\omega}$ et
pour les mêmes rapports entre $\mathcal{A}$ et $B$, dans les équations
qui conviendront, \struck{aux cas semblables pour} la plaque quarrée.

\struck{Demême aussi que} On peut conclure demême que les mouvemens
d'une surface triangulaire isocèle rectangle dont l'hypothénuse
seroit appuyée, un des côtés adjacens al'angle droit libre en vertu des
conditions $\frac{dz}{dx}=0$ et $\frac{d\left\{\frac{ddz}{dx^2}+\frac{ddz}{dy^2}\right\}}{dx}=0$ et l'autre côté
\note{En tête, au milieu, « (59) », barré ; en haut à droite, à l'encre,
« 229 ». Le feuillet est collé bas sur sa feuille de montage, comme ceux
des vues 461 et 463 : le premier tiers de la vue est blanc. La première
ligne achève la phrase de la vue 464 (« qu'aux seules surfaces
rectangulaires / car ces surfaces… ») : la pagination de Sophie Germain
passe de (58) à (59) par-dessus le fragment de la vue 465. « toujours »
est écrit au-dessus de « puisse être » ; « , par exemple, » au-dessus de
la fin de la première formule. Dans cette formule, le $d$ devant
l'accolade n'a pas de dénominateur ; à la dernière ligne, la même
condition est écrite avec « $dx$ » sous un trait de fraction : les deux
sont gardées telles qu'elles sont écrites. « qui se meut » est barré
après « quarrée ». « $x$ », « $y$ », $\omega$ sont soulignés. « nœuds »
est écrit « nauds », l'œ ouvert, comme aux vues 350–356. « Ainsi s'il »
est barré et « Si donc il » écrit au-dessus. Le « n » de « nodales »
est repassé. Après « au sens d' », le nom, dont l'E est tracé comme un L, est lu
« Euler », comme aux vues 419–420. « aux cas semblables pour » et « Demême aussi que »
sont barrés. La phrase se poursuit à la vue suivante.}

\page{468}

adjacent au même angle libre en vertu des conditions
\struck{$\frac{dz}{dy}=0$ et} $\frac{d\left\{\frac{ddz}{dx^2}+\frac{ddz}{dy^2}\right\}}{dy}=0$ et $\frac{ddz}{dx^2}+\frac{ddz}{dy^2}=0$, seront
compris également dans les équations qui appartiennent, pour
les differens cas respectifs, au mouvement de la plaque
quarrée; \struck{Et demême} \add{Il faut conclure} encore \add{que les mouvemens} \struck{qu'} d'une pareille surface
triangulaire dont l'hypothénuse et un des côtés adjacent a
l'angle droit seroient libres au sens d'Euler et l'autre appuyé,
seroient aussi compris dans les équations qui déterminent
les mouvemens dela plaque quarrée. Il faudra seulement
\struck{observer}, pour transporter à ces differens triangles l'analyse
qui appartient ala plaque quarrée, \add{avoir soin} de changer, dans les
coëfficiens, $\underline{\omega}$ en $\frac{1}{2}\omega$ \&c. suivant \struck{qu'il s'agit d'un triangle
dont les côtés sont} le rapport qui existe entre le côté
du triangle que l'on considère et celui dela plaque
entière, \struck{c'est} J'ai déjà fait une remarque semblable en
parlant de la simple lame et des rapports qui existent entre
les differens cas qu'Euler a calculés a part.

Je ne m'arreterai pas à appuyer et à developper les observations
que je viens de \struck{faire} \add{présenter} j'ai voulu seulement les indiquer \struck{pour \ill{}} \add{afin de donner}
\struck{d'} \add{un} exemple d\struck{e l'\uncertain{avantage}}\add{u} parti que l'on peut tirer dela considération
des points et des lignes de limites pour déterminer, au moyen
d'équations déjà connues le mouvement de surfaces qui, au
premier aperçu, ne sembloit \add{pas} devoir y être compris.
\note{En tête, au milieu, « (60) », non barré ; aucun nombre en haut à
droite : c'est le dos du feuillet de la vue 467. Au-dessus du feuillet,
sur la feuille de montage, paraît à l'envers, par transparence, le
fragment de la vue 465 (« 228 »). La première ligne achève la phrase de
la vue 467 (« et l'autre côté / adjacent »). Au début de la formule,
« $\frac{dz}{dy}=0$ et » est barré de traits épais. Après « quarrée ; »,
« Et demême » est barré ; « Il faut conclure » est écrit au-dessus,
et « que les mouvemens » au-dessus de « encore » ; le « qu' » devant
« d'une » est barré. « observer » est barré en tête de ligne, et « avoir
soin » écrit au-dessus de la ligne suivante, après « quarrée, ». « ces »
est écrit sur « les » devant « differens triangles ». « qu'il s'agit d'un
triangle / dont les côtés sont » est barré sur deux lignes ; « c'est » est
barré devant « J'ai ». Le nom d'Euler a son E tracé comme un L. Dans le
dernier alinéa, écrit un peu en retrait : « faire » est barré et
« présenter » écrit au-dessus ; « pour » et le mot suivant sont barrés en
fin de ligne, et « afin de donner » écrit au-dessus ; en tête de la ligne
suivante, « d' » est barré et « un » écrit au-dessus ; « du » est formé
sur « de l' », et le mot barré qui suit est lu avec doute « avantage » ;
« pas » est écrit au-dessus de « devoir ». Le texte s'arrête ici, au
milieu de la page, sur un point : la fin d'une section. Le reste du
feuillet est blanc.}

\page{469}

Je ne puis terminer cet article sans rendre \struck{\uncertain{homage}} \add{hommage} à l'étonnante
sagacité de M\textsuperscript{r} Chladni qui, sans avoir connu la théorie ni
parconséquent l'étendue de l'intégrale (C), s'est pourtant décidé
à regarder les figures, en apparence si bizarres, qui y sont
comprises, comme parfaitement équivalentes à de simples lignes
droites parallelles entre elles: j'avoue même que j'avois d'abord
envisagé le rapport que l'habile physicien indique entre
les figures dont il s'agit comme une simple classification
fondée uniquement sur la possibilité d'obtenir successivement
plusieurs de ces figures, sans que le son varie. J'étois
alors bien loin de prévoir qu'un simple changement de
rapport entre des constantes, pût rendre compte dela variété
des figures, tandis que la constance de la valeur de l'angle
$\underline{\omega}$ dans celles de ces figures que M\textsuperscript{r} Chladni avoit jugées
equivalentes, expliqueroit naturellement l'uniformité des sons
qu'il avoit observés.

N.o 13 \emph{Surfaces périodiques}. Je passe à présent à l'application
des intégrales données N.o 4 et 5. Leur forme générale est exprimée
par la formule $z = \ldots\ldots\,\text{sin}\left(\pi.\frac{nx}{A}\right)\text{sin}\left(\pi.\frac{my}{A}\right)$; et cette formule
montre que, si en l'appliquant à une surface quarrée, on prend
l'origine à un des angles de la surface, lamême surface sera
appuyée dans tous les points de son contour; car aux extrêmités,
c'est-adire, lorsque \struck{$x=A$ et $x=0$ ou $y=A$ et $y=0$} l'on aura une
des quatre conditions $x=A$, $x=0$, $y=A$, $y=0$
\note{En tête, au milieu, « (61) », barré ; en haut à droite, à l'encre,
« 230 ». Le feuillet couvre les deux tiers supérieurs de la vue ;
au-dessous, la feuille de montage. Un paragraphe nouveau ouvre la page ;
la page (60), vue 468, s'était arrêtée sur un point au milieu du
feuillet. Devant « à l'étonnante », un mot barré, lu avec doute
« homage », est récrit au-dessus « hommage ». Le numéro du paragraphe
se lit « N.o 13 », le 3 à queue courte ; les N.o 11 et 12 ne sont pas
dans ce lot. « Surfaces périodiques » est souligné. « N.o 4 et 5 » :
le 4 est tracé comme un « h ». Dans la formule, les points de
suspension sont les siens, pour un facteur laissé en blanc ; « sin »
est écrit avec le s long. Une croix sous la ligne, entre « c'est-adire, »
et « lorsque », ne renvoie à rien de visible. « $x=A$ et $x=0$ ou
$y=A$ et $y=0$ » est barré, et les quatre conditions sont récrites à la
ligne suivante, séparées par des virgules. La phrase se poursuit à la
vue suivante.}

\page{470}

les deux équations $z=0$ et $\frac{ddz}{dx^2}+\frac{ddz}{dy^2}=0$ seront satisfaites

J'ai déjà dit que le mot \emph{appuyé} ne devoit être pris
que dans un sens analytique et comme equivalent aux
deux conditions $z=0$ et $\frac{ddz}{dx^2}+\frac{ddz}{dy^2}=0$. Ainsi je l'emploirai
dorénavant sans revenir davantage sur sa signification; et je
nommerai \emph{ligne\struck{s} d'appui\struck{s}} la suite des points qui satisferont
aux conditions qui \struck{en} fixent la valeur de ce mot.

Cela posé, une periode complette\add{, tant dans le sens des $x$ que dans celui des $y$,} sera comprise entre quatre
lignes d'appui: ainsi, en prenant toujours l'origine à un des
angles de la plaque\add{, lorsque l'on aura $m=n=1$,} la courbure de cette surface sera
composée d'une seule periode \struck{lorsque l'on aura $m=n=1$};
mais, \add{dans un grand nombre de cas,} cette période sera partagée en quatre quarts de
périodes par deux lignes de limites analytiques qui se couperont
a angle droit au centre de la surface, et satisferont aux
conditions $\frac{dz}{dx}=0$ et $\frac{d\left\{\frac{ddz}{dx^2}+\frac{ddz}{dy^2}\right\}}{dx}=0$ ou $\frac{dz}{dy}$ et $\frac{d\left\{\frac{ddz}{dx^2}+\frac{ddz}{dy^2}\right\}}{dy}=$
qui pourroient également appartenir aux extrêmités libres.

Il suit delà que l'équation qui, \add{dans les cas dont j'ai parlé,} \struck{sous des conditions données}
appartient àla surface quarrée dont les quatre côtés sont
appuyés, appartiendroit également, et sous les mêmes conditions,
à une surface quarrée dont les quatre côtés seroient libres.
En effet les deux cas ne different que par la disposition
des portions de periodes dans lesquelles se trouvent les points
de contour dela surface; desorte que, par exemple,
lorsqu'il s'agit d'une seule période, la position des quatre
\note{En tête, au milieu, « (62) », non barré ; aucun nombre en haut à
droite : c'est le dos du feuillet de la vue 469. La première ligne
achève la phrase de la vue 469 (« l'on aura une des quatre conditions
… / les deux équations … seront satisfaites ») ; elle finit sans
point. « appuyé » est souligné. Dans « lignes d'appuis », souligné, le
« s » final des deux mots est barré. « en » est barré devant « fixent ».
« tant dans le sens des x que dans celui des y » est écrit au-dessus de
« complette, sera » ; « , lorsque l'on aura $m=n=1$, » au-dessus de
« angles de la plaque », et la même condition est barrée à la ligne
suivante ; « dans un grand nombre de cas, » est écrit au-dessus de
« mais, cette ». Dans la dernière ligne de formules, « $\frac{dz}{dy}$ »
est écrit sans « $=0$ », et la seconde condition finit sur « = » au
bord du feuillet, sans le zéro : l'une et l'autre sont gardées ainsi.
« dans les cas dont j'ai parlé, » est écrit au-dessus de « sous des
conditions données », barré. La phrase se poursuit à la vue suivante.}

\page{471}

quarts de période\struck{s}, dans un des cas, est exactement le renversement
dela disposition qui a lieu dans l'autre; et qu'en supposant
ces quatre portions de période\struck{s} isolés, elles doivent être rapprochés
par leurs côtés libres, pour former la surface dont les quatre côtés
sont appuyés; et par leurs côtés appuyés, pour former
la surface dont les quatre côtés sont libres. J'ai déjà remarqué
un rapport semblable à celui que je viens d'indiquer, entre
la corde tendue et la flûte ouverte aux deux bouts (V N.o 7).

\struck{Après avoir \uncertain{présenté} \uncertain{l'observation} précédente qui \uncertain{s'applique} a
toutes les formules périodiques je vais examiner en particulier
les conséquences de l'intégrale (\ill{}).}

L'équation générale $z = \ldots\ldots\,\text{sin}.\left(\pi.\frac{nx}{A}\right)\text{sin}\left(\pi.\frac{my}{A}\right)$ montre
encore qu'il y aura toujours sur la plaque quarrée àla
quelle elle sera appliquée, $\underline{n}$ lignes nodales parallelles à l'axe
des $\underline{y}$, c'est adire à un des côtés de la plaque, et $\underline{m}$ lignes
nodales parallelles à l'axe des $\underline{x}$, et parconséquent à l'autre
côté de la plaque.+ Elle montre aussi que \struck{ces lignes seront
toujours a égales distances} celles de ces lignes qui seront parallelles
entre elles, se trouveront toujours à égale distance, les unes des
autres; que cette distance sera double de celle entre la dernière
des mêmes lignes et le côté libre auquel elle \struck{est} \add{sera} parallelle;
que chaque intervalle de deux lignes nodales parallelles \struck{est} \add{sera}
partagé par une ligne de limite analytique qui \uncertain{satisfera}
à l'un ou l'autre des couples d'équations $\frac{dz}{dx}=0$ et $d.\left\{\frac{ddz}{dx^2}+\frac{ddz}{dy^2}\right\}=0$
\note{En tête, au milieu, « (63) », barré ; en haut à droite, à l'encre,
« 231 ». Le feuillet est collé bas sur sa feuille de montage : le premier
tiers de la vue est blanc. La première ligne achève la phrase de la
vue 470 (« la position des quatre / quarts de période »). Dans « périodes »,
aux deux premières lignes, le « s » final est barré d'une petite croix.
« (V N.o 7) » renvoie à un paragraphe qui n'est pas dans ce lot. Les
trois lignes données ici en \emph{struck} sont barrées lettre par lettre
de petites croix, comme aux vues 314, 404 et 405 ; les mots en sont lus
sous les croix, plusieurs avec doute, et l'étiquette de l'intégrale, entre
parenthèses à la fin, ne se lit pas. Dans les deux formules, « sin » est
écrit avec le s long, et les points de suspension sont les siens.
« $n$ », « $y$ », « $m$ », « $x$ » sont soulignés. Une croix suit « de la
plaque. » ; elle n'appelle rien de visible sur ce feuillet. « ces lignes
seront / toujours a égales distances » est barré sur deux lignes ; « est »
est barré deux fois et « sera » écrit au-dessus. La fin de « satisfera »
est surchargée et barrée. Dans la dernière formule, « et » est écrit
au-dessus du signe, entre les deux équations. La phrase se poursuit à la
vue suivante.}

\page{472}

$\frac{dz}{dy}=0$ \quad $\frac{d\left\{\frac{ddz}{dx^2}+\frac{ddz}{dy^2}\right\}}{dy}=0$; et enfin que tout point
pour lequel on aura àlafois $\text{sin}\left(\pi.\frac{nx}{A}\right)=0$ et $\text{sin}\left(\pi.\frac{my}{A}\right)=0$,
c'est adire tout point d'intersection des lignes nodales
droites qui terminent les periodes, pourra être pris pour
origine, puisque, conformement à cequej'ai dit N.o 9,
il satisfait aux quatre conditions $z=0$, $\frac{dz}{dx}=0$,
$\frac{ddz}{dx^2}+\frac{ddz}{dy^2}=0$ et $\frac{d\left\{\frac{ddz}{dx^2}+\frac{ddz}{dy^2}\right\}}{dx}=0$.

Après avoir exposé les conséquences générales dela périodicité
des surfaces, je vais examiner en particulier celles
qui resultent de l'intégrale (D). N.o 4

On voit d'abord que les surfaces quarrées dont les
mouvemens seront compris dans cette intégrale, ne pourront
présenter d'autres lignes nodales que \struck{celle} celles
qui terminent les périodes; car aucune autre supposition
que $\text{sin}\left(\pi.\frac{nx}{A}\right)=0$ ou $\text{sin}\left(\pi.\frac{my}{A}\right)=0$ ne pourra
alors rendre $z=0$.

Parmi\struck{s} les cas observés par M\textsuperscript{r} Chladni, ce sont ceux
auxquels appartiennent les figures 63, 66a, 68a,
71a, 74a, 75a, 80a, 81a et 82 (V. T. d'acoustique)
qui sont renfermés dans l'équation (D).

Pour la figure 63; $m=n=1$., \struck{et, conformément à l'experience,}
il ne doit \add{donc} y avoir+ que deux lignes nodales parallelles:
\marginal{+, conformément à l'experience,}
\note{En tête, au milieu, « (64) », non barré ; aucun nombre en haut à
droite : c'est le dos du feuillet de la vue 471. La première ligne
achève la phrase de la vue 471 (« à l'un ou l'autre des couples
d'équations … / $\frac{dz}{dy}=0$ … ») ; entre les deux équations du
second couple, aucun « et » n'est écrit. « sin » est écrit avec le s
long. « N.o 9 » renvoie au paragraphe sur les plaques, à la vue 433
(lot 22). L'étiquette de l'intégrale, deux fois, est une capitale
empâtée entre parenthèses, lue « D » : c'est l'intégrale
$z = \mathcal{A}\,\text{sin}\frac{\pi nx}{A}\,\text{sin}\frac{\pi my}{A}$ du N.o 4
(vue 406, lot 21), et le numéro qui suit « N.o », surchargé, se lit 4. « celle » est barré et récrit « celles ». Le « s »
de « Parmis » est barré ; « ce » est écrit sur un mot plus court. Les
numéros de figures sont ceux du traité d'acoustique de Chladni (« V. T.
d'acoustique »). À la dernière phrase, « et, conformément à
l'experience, » est barré sur la ligne, et la même incise est récrite
dans la marge de gauche, appelée par une croix après « y avoir » ;
« donc » est écrit au-dessus de « doit ». La phrase se poursuit à la vue
suivante.}

\page{473}

l'une à l'axe des $\underline{x}$, l'autre à celui des $\underline{y}$; et ces lignes doivent
se couper au centre dela surface dans un point d'origine.
\struck{qui est dont \uncertain{cette} être le \ill{}}
\struck{Le son est parconséquent} \add{le son doit être} proportionnel au nombre
$m^2+n^2=2$ (V. N.o 4).

Pour la figure 66a $m=2$, $n=1$; \struck{et} conformement a l'experience
il \struck{ne} doit \add{donc} y avoir deux lignes nodales parallelles à l'axe des
$\underline{x}$ et une seule parallelle à l'axe des $\underline{y}$.
\struck{Parconséquent} le son doit être proportionnel au nombre
$m^2+n^2=5$.

Pour la figure 74a $m=4$, $n=1$; conformement à l'experience
il doit donc y avoir 4 lignes nodales parallelles à l'axe
des $x$ et une seule parallelle à l'axe des $\underline{y}$.
Le son doit être proportionnel au nombre $m^2+n^2=17$.

Les trois valeurs précédentes de $\underline{n}$ étant 1, elles appartiennent
a un cas \struck{des vibrations} que M\textsuperscript{r} Chladni a cru devoir
distinguer, et qu'il a même proposé sous le nom de
\emph{vibrations tournantes}, comme pouvant aider à déterminer
le reste des mouvemens dont les plaques élastiques
sont susceptibles (V. T. d'acoustique N.o 84 $\ldots$ 87 et N.o 98)

D'après la théorie, il est clair que les \emph{vibrations tournantes}
ne diffèrent des autres vibrations \struck{des plaques} comprises
\note{En tête, au milieu, « (65) », barré ; en haut à droite, à l'encre,
« 232 ». Le feuillet couvre les deux tiers supérieurs de la vue. La
première ligne achève la phrase de la vue 472 (« deux lignes nodales
parallelles: / l'une à l'axe des x »). Après « celui des y », le
point-virgule est surchargé par « et ». La troisième ligne est barrée
entière ; seuls « qui est dont » et « être le » s'y lisent, le reste
avec doute ou pas du tout. « Le son est » et « parconséquent » sont barrés,
et « le son doit être » écrit au-dessus. « (V. N.o 4) » renvoie à un
paragraphe qui n'est pas dans ce lot. « et » est barré devant
« conformement », « ne » devant « doit » ; « donc » est écrit au-dessus
de « doit y avoir ». « Parconséquent » est barré en tête de ligne. Le 4
de « 74a », « $m=4$ », « 4 lignes » et « N.o 84 » est tracé comme un
« h ». « vibrations tournantes » est souligné
les deux fois. Les points de suspension de « N.o 84 $\ldots$ 87 » sont
les siens. La phrase se poursuit à la vue suivante.}

\page{474}

dans l'intégrale (D) qu'en ceque \struck{elles ne sont composées} les surfaces
auxquelles elles appartiennent n'ont qu'une seule periode dans le
sens des $x$, période qui pour les plaques dont les \struck{extrêmités} \add{côtés}
sont libres, est composée de deux demies periodes reunies
dans la ligne de repos parallelle à l'axe des $\underline{y}$.
Cette difference n'est pas analytiquement importante;
cependant j'aurai occasion\add{ de remarquer}, en traitant dela comparaison
entre les sons \struck{d} indiqués par la théorie et les sons donnés
par l'experience, que le cas des \emph{vibrations tournantes}
présente une singularité digne d'attention.

Pour la fig. 68.a $m=2$ $n=2$; Conformement à l'experience
il doit donc y avoir 2 lignes nodales parallelles à l'axe
des $\underline{x}$, et deux lignes nodales parallelles à l'axe des $\underline{y}$.
Le son doit être proportionnel au nombre $m^2+n^2=8$.

Pour la fig. 71 a $m=2$, $n=3$; Conformement àl'expérience
il doit y avoir 2 lignes nodales parallelles à l'axe des $\underline{x}$
et 3 lignes nodales parallelles à l'axe des $\underline{y}$
Le son doit être proportionnel aunombre $m^2+n^2=13$.

Pour la fig. 75a $m=3$ $n=3$. Conformement à l'expérience,
il doit donc y avoir 3 lignes nodales parallelles à chacun
des axes $\underline{x}$ et $\underline{y}$.
Le son doit être proportionnel au nombre $m^2+n^2=18$
\note{En tête, au milieu, « (66) », non barré ; aucun nombre en haut à
droite : c'est le dos du feuillet de la vue 473. La première ligne
achève la phrase de la vue 473 (« ne diffèrent des autres vibrations
comprises / dans l'intégrale (D) »). L'étiquette de l'intégrale est la
même capitale empâtée qu'à la vue 472, lue « D ». « elles ne
sont » et « composées » sont barrés d'un même trait. « extrêmités » est
barré et « côtés » écrit au-dessus. « de remarquer, » est écrit au-dessus
de « occasion, en » ; après « les sons », une lettre commencée est barrée.
« vibrations tournantes » est souligné. Les numéros de figures sont ceux
du traité de Chladni ; dans « 68.a », « 71 a », « 75a », les chiffres
sont nets. Dans « $m=3$ $n=3$ », le 3 est à queue courte. Le texte
s'arrête ici, au bas du feuillet, sans point ; la page suivante (vue 475)
reprend sur un alinéa nouveau.}

\page{475}

Pour la fig. 80a $m=3$, $n=4$. Conformement àl'experience, il doit donc y avoir
3 lignes nodales parallelles à l'axe des $\underline{x}$, et quatre ligne nodales parallelles
à l'axe des $\underline{y}$.
Le son doit être proportionnel au nombre $m^2+n^2=25$.

Pour la fig. 81a. $m=2$ $n=5$. Conformement à l'experience, il doit donc
y avoir 2 lignes nodales parallelles à l'axe des $\underline{x}$, et 5 lignes nodales
parallelles à l'axe des $\underline{y}$.
Le son doit être proportionnel au nombre $m^2+n^2=29$.

Enfin pour la fig. 82. $m=4$. $n=4$. Conformement àl'experience, il doit
donc y avoir 4 lignes nodales parallelles à chacun des axes $\underline{x}$ et $\underline{y}$.
Le son doit être proportionnel au nombre $m^2+n^2=32$.

Je vais m'occuper àprésent des valeurs de $\underline{z}$ qui donnent lieu à
l'existence de figures nodales plus compliquées que celles qui sont
\struck{renfermés dans} appartiennent ala formule (D)

N.o 14 Les intégrales (E) et (F) N.o 4 contiennent, outre le facteur
dépendant dela periodicité, un autre facteur, sous forme de
série; \struck{\uncertain{duquel}} lequel détermine la figure nodale comprise dans
chaque periode ou partie\struck{s} de periode. Dans ce facteur,
les constantes $a'$, $b'$, $A'$ et $B'$ sont independantes entr'elles,
desorte que, pourvu qu'on ne fasse pas àlafois $a'=0$, $b'=0$ ou
$A'=0$, $B'=0$, suppositions qui donneroient $z=0$., on peut imaginer
entre ces constantes \uncertain{telle} relation qu'il plaira de supposer.
Il est clair que la figure nodale comprise dans chaque
\note{En tête, au milieu, « (67) », barré ; en haut à droite, à l'encre,
« 233 ». Le feuillet est collé bas sur sa feuille de montage : le premier
tiers de la vue est blanc. La page commence sur un alinéa nouveau, à la
suite de la série de figures de la vue 474. Dans « 80a », le 8 est
surchargé. « ligne nodales » : le « s » manque, ou n'a pas été tracé.
Dans la dernière égalité, $m^2+n^2=32$, l'exposant de $m$ est à peine
formé. « renfermés dans » est barré. L'étiquette de la formule est la
capitale empâtée des vues 472 et 474, lue « D ». Le paragraphe
est numéroté « N.o 14 », à la suite du N.o 13 de la vue 469 ; le 4 est
tracé comme un « h », comme dans « N.o 4 » à la même ligne. « duquel »,
lu avec doute, est surchargé par « lequel ». Le « s » de « parties »
est barré. Dans « donneroient », la fin est surchargée. Après « ces
constantes », un mot surchargé plusieurs fois se lit avec doute « telle ».
$A'$ et $B'$ sont des capitales droites, distinctes de la capitale cursive
de la vue 461. La phrase se poursuit à la vue suivante.}

\page{476}

période, dépendra dela relation établie entre les constantes
$a'$, $b'$, $A'$ et $B'$: mais les sons seront aucontraire indépendans
de cette relation. Ainsi on voit qu'ici, comme dans les cas
renfermés dans l'intégrale (C), plusieurs figures différentes
peuvent rendre un seul et même son.

Pour les cas qui appartiennent à l'intégrale (E), le
son est proportionnel au nombre $(\Pi-1)(m^2+n^2)$ (V. N.o 4).
ainsi, en \struck{prenant $\Pi=1-p^2$, et} observant que\struck{, en general}
dans les formules qui expriment les sons, tant dans la présente
intégrale (E) que dans les suivantes, $(\Pi-1)$ \struck{\uncertain{vient}} \add{vient} de
$\sqrt{(\Pi-1)}^2$, \struck{et qu'on} \add{et} \add{parconséquent} peut changer le signe de \struck{$(\Pi-1)$} \add{ce nombre} on voit
qu'en prenant $\Pi=1-p^2$, le son devient proportionnel
au nombre $p^2(m^2+n^2)$ \struck{et} qui \struck{ce nombre} est dela même
forme que celui \struck{auquel} \add{qui exprime} les sons \struck{qui} appartenant
aux manières de vibrer renfermés dans l'intégrale
(D)

Demême, pour les cas qui appartiennent à l'intégrale
(F), le nombre $m^2(\Pi+\Pi'-2)$ auquel le son est proportionnel,
devient\add{,} \struck{en prenant $\Pi=1-p^2$, $\Pi'=1-q^2$ et changeant
les signes,} $m^2(p^2+q^2)$, en prenant $\Pi=1-p^2$, $\Pi'=1-q^2$ et
changeant les signes. On voit donc que le son est
alors proportionnel à un nombre de même forme
que \struck{\uncertain{celui}} pour les cas renfermés dans l'intégrale (D)
\note{En tête, au milieu, « (68) », non barré ; aucun nombre en haut à
droite : c'est le dos du feuillet de la vue 475. La première ligne
achève la phrase de la vue 475 (« la figure nodale comprise dans chaque
/ période, dépendra »). La lettre des formules des sons est un grand
$\Pi$ aux jambages hauts, rendu $\Pi$ et $\Pi'$ comme au lot 21 (vues
407–413), où l'intégrale (E) donne le son proportionnel à
« $\Pi-1$ » ; elle se distingue du $\pi$ des sinus des vues 469–472.
« (V. N.o 4) » : le 4 est tracé comme un « h ». « prenant $\Pi=1-p^2$,
et » est barré ; après « que », « , en general » est barré. Après
« $(\Pi-1)$ », un mot est surchargé et barré, et « vient » est écrit
au-dessus. Devant « peut », « et qu'on » est barré, et « et » puis
« parconséquent » sont écrits au-dessus. « $(\Pi-1)$ » est barré après
« le signe de », et « ce nombre » écrit au-dessus. « et » et « ce
nombre » sont barrés dans la ligne « au nombre $p^2(m^2+n^2)$ … » ;
« auquel » est barré et « qui exprime » écrit au-dessus ; « qui » est
barré devant « appartenant », dont la fin est surchargée.
L'étiquette de l'intégrale, deux fois à la fin des alinéas, est la
capitale empâtée des vues 472–475, lue « D ». Devant
« devient », un signe d'insertion en forme de virgule est écrit sous la
ligne. « en prenant … changeant / les signes, » est barré sur deux lignes
et récrit après « $m^2(p^2+q^2)$ ». Devant « pour les cas », un mot court,
lu avec doute « celui », est barré. Le texte s'arrête ici, au bas du
feuillet, sans point ; la page suivante reprend sur un alinéa nouveau.}

\page{477}

Si dans l'intégrale (E) on \uncertain{fait} $\Pi=1-p^2$, et que l'on \struck{prenne} fasse
ala fois $b'=0$ et $B'=0$, ou $a'=0$ et $A'=0$\add{, suivant que $p$ est pair ou impair}, nonseulement
les sons appartenans aux manières devibrer renfermés dans
cette intégrale seront exprimés par les mêmes nombres
que \struck{dans} \add{lorsqu'il s'agit de} l'intégrale (D), mais encore ces deux formules
deviendront \struck{identiques}. \struck{équivalentes}. identiques

En effet \add{si dans} le facteur $\text{sin}.\left(\frac{\pi nx}{A}\right)\left\{\text{cos}\left(\frac{\pi nx}{A}\right) + \frac{4-p^2}{2.3}\text{cos}^3\left(\frac{\pi nx}{A}\right) + \frac{(4-p^2)(16-p^2)}{2.3.4.5}\text{cos}^5\left(\frac{\pi nx}{A}\right)\right.$
$\left. + \&c\right\}$
on change sin. en cos. et \emph{vice versa} c'est adire, si \struck{on} \add{l'on} met $\underline{\frac{1}{2}A-x}$
àla place de $\underline{x}$, on trouvera, après avoir multiplié par la
constante $p$, que le \struck{facteur, \uncertain{soit}} même facteur ne sera autre chose
que l'expression de $\text{sin}\left(\pi.\frac{n.px}{A}\right)$ donnée par la formule P. 143
\emph{des leçons sur le calcul des fonctions} \add{\uncertain{éd.}} \emph{in.8 an 1806}.

et demême si dans le facteur
$\text{sin}\left(\frac{\pi nx}{A}\right)\left\{\frac{2}{1-p^2} + \text{cos}^2\left(\frac{\pi nx}{A}\right) + \frac{9-p^2}{3.4}\text{cos}^4\left(\frac{\pi nx}{A}\right) + \&c\right\}$ on change
$x$ en $\underline{\frac{1}{2}A-x}$ et que l'on multiplie par la constante $\frac{1-p^2}{2}$ on
trouvera qu'il n'est autre chose que l'expression de $\text{cos}\left(\pi.\frac{npx}{A}\right)$
donnée P. 143 \emph{des leçons}.

On aura donc alors $z = \ldots\ldots\,\text{sin}.\left(\pi.\frac{npx}{A}\right)\text{sin}\left(\pi.\frac{mpy}{A}\right)$ ou cequi revient
au même $z = \ldots\ldots\,\text{cos}.\left(\pi.\frac{npx}{A}\right)\text{cos}\left(\pi.\frac{mpy}{A}\right)$.

A l'égard de l'intégrale (E), on trouvera également en prenant
$\Pi=1-p^2$, $\Pi'=1-q^2$, et fesant àlafois $b'=0$ et $B'=0$, ou $a'=0$ et
$A'=0$, qu'elle se réduit à l'une ou l'autre des formes
équivalentes. $z = \ldots\,\text{sin}\left(\pi.\frac{mpx}{A}\right)\text{sin}\left(\pi.\frac{mqy}{A}\right)$ ou $z = \ldots\,\text{cos}\left(\pi.\frac{mpx}{A}\right)\text{cos}.\left(\pi.\frac{mqy}{A}\right)$.
\note{En tête, au milieu, « (69) », barré ; en haut à droite, à l'encre,
« 234 », le 4 empâté. La page ouvre un alinéa nouveau, à la suite de la
vue 476. Après « on », un mot est surchargé ; la lecture « fait » est
offerte avec doute. « prenne » est barré et « fasse » écrit à la suite.
« suivant que $p$ est pair ou impair » est écrit au-dessus de « $A'=0$,
nonseulement ». « dans » est barré et « lorsqu'il s'agit de » écrit
au-dessus ; l'étiquette de l'intégrale est la capitale empâtée des vues
472–476, lue « D ». « identiques » est écrit une première fois
et barré, puis « équivalentes » est barré et « identiques » récrit à la
suite. La lettre des coefficients est le même grand $\Pi$ qu'à la vue
476 ; celle des sinus et cosinus, le $\pi$ des vues 469–472. Le « 4 »
de « $4-p^2$ » est tracé comme un « h » ; le premier chiffre du second
facteur, « $16-p^2$ », se lit aussi 9 (« 96 ») et n'est pas sûr. « si
dans » est écrit au-dessus de « En effet ». « on » est surchargé par
« l'on ». « vice versa », « $\frac{1}{2}A-x$ » (deux fois) et « $x$ »
sont soulignés ; « des leçons sur le calcul des fonctions in.8 an 1806 »
est souligné, et « des leçons » à la ligne « donnée P. 143 ». Au-dessus de
« in.8 », deux lettres suivies d'un point sont lues avec doute « éd. ».
Le « P » de « P. 143 » est empâté la première fois, et tracé comme un 8
la seconde. « facteur, » et un mot court lu avec doute « soit » sont
barrés avant « même facteur ». Dans « cos.(…)cos », à la ligne « au même »,
le premier « cos » est écrit sur un « sin ». Dans la dernière formule, le
second « cos » est écrit sur un autre mot, et « équivalentes » a sa fin
surchargée. Les points de suspension des formules sont les siens. Le
dernier alinéa porte « l'intégrale (E) », nettement ; les deux
constantes $\Pi$, $\Pi'$ et le nombre $m^2(\Pi+\Pi'-2)$ de la vue 476
sont ceux de l'intégrale (F) : le texte est gardé tel qu'il est écrit. Le
texte s'arrête ici, aux deux tiers de la vue, sur un point.}

\page{478}

Pour tout autre supposition que celles qui ramenent à l'intégrale
(D), les séries contenus dans les intégrales (E) et (F) ne se
terminent pas; \struck{et} il faudroit \add{donc} calculer àpart chaque
cas, pour déterminer la figure qui lui convient. \struck{Mais}
Cependant les valeurs de $\Pi$ et $\Pi'$ feront connoître tout
de suite le nombre des lignes nodales qui devront
composer ces figures, \struck{Mais} je ne m'arreterai pas
à ces considérations. \struck{et je me bornerai, après avoir
présenté quelques unes des conséquences générales dela
forme des intégrales (E) et (F)} Il me suffira d'avoir
présenté quelques unes des conséquences générales dela
forme des intégrales (E) et (F), et je passerai à
l'examen de l'intégrale (G) qui \struck{présente est} susceptible
, dans certains cas, d'une application indépendante de
tout calcul.

N.o 15. Si dans l'équation (G) on fait $a''=0$, $b''=0$, $A''=0$
et $B''=0$, elle se réduira à l'équation (F). Ce cas est
le seul où la même équation (G) puisse devenir
identique à l'intégrale (D).

Si aucontraire, dans l'intégrale (G) on fait $a'=0$, $b'=0$
$A'=0$ et $B'=0$, $z''$ sera la somme de deux fonctions
séparées des variables $\underline{x}$ et $\underline{y}$; \struck{desorte que} \add{et} tandis que
\note{En tête, au milieu, « (70) », non barré ; aucun nombre en haut à
droite : c'est le dos du feuillet de la vue 477. La page ouvre un alinéa
nouveau. L'étiquette « (D) », deux fois, est la capitale empâtée des
vues 472–477. « et » est barré devant « il faudroit » ;
« donc » est écrit au-dessus. « Mais » est barré en fin de ligne, puis
une seconde fois après « ces figures, ». Les trois lignes « et je me
bornerai … (E) et (F) » sont barrées et récrites à la suite, sous une
autre tournure. « présente est » est barré devant « susceptible » ; la
virgule qui suit est reportée en tête de la ligne suivante. Le
paragraphe est numéroté « N.o 15 », le 5 à queue descendante ; il suit
le N.o 14 de la vue 475. Les capitales $A''$, $B''$, $A'$, $B'$ ont
les accents écrits au-dessus d'une capitale droite. La lettre avant
« sera la somme » est un $z$ muni de deux traits, lu $z''$. « $x$ » et
« $y$ » sont soulignés. « desorte que » est barré et « et » écrit
au-dessus. La phrase se poursuit à la vue suivante.}

\page{479}

$z'$ satisfera \struck{a la condit.} l'équation $\frac{ddz'}{dx^2} = \frac{n^2}{m^2}.\frac{ddz'}{dy^2}$, $z''$ remplira la
condition $\frac{ddz''}{dx\,dy}=0$.

Si on veut appliquer l'intégrale (G) ainsi réduite, au cas où les côtés
dela plaque quarrée sont libres en vertu des conditions
$\frac{dz}{dx}=0$ et $\frac{d\left\{\frac{ddz}{dx^2}+\frac{ddz}{dy^2}\right\}}{dx}=0$ ou $\frac{dz}{dy}=0$ et $\frac{d\left\{\frac{ddz}{dx^2}+\frac{ddz}{dy^2}\right\}}{dy}=0$
on trouvera qu'il faut encore faire $a''=0$ et $A''=0$.
\add{car sans cette condition, elle ne conviendroit qu'aux plaques dont les côtés seroient appuyés.}
Cela posé, en prenant $\Pi''=1-9$, l'équation (G) deviendra
\[ z = \ldots\,\text{sin}\left(\frac{\pi mx}{A}\right).\,\text{sin}\left(\frac{\pi my}{A}\right)\left\{b\,\text{cos}^2\left(\frac{\pi mx}{A}\right) + B\,\text{cos}^2\left(\frac{\pi my}{A}\right) - \frac{(b+B)}{4}\right\} \]

Parmi les figures observées par M\textsuperscript{r} Chladni celles qui doivent
être rapportées à cette intégrale sont les fig. 69, 70,
87a, 87b, 88a et 88b.

Pour la fig. 69. $m=1$, $b=-B$ conformement àl'experience, il doit y avoir
outre les deux lignes nodales dependantes du facteur periodique deux autres
lignes nodales diagonales dépendantes du facteur $\text{cos}^2\left(\frac{\pi x}{A}\right) - \text{cos}^2\left(\frac{\pi y}{A}\right) = 0$.

Pour la fig. 70. $m=1$; mais il existe entre $b$ et $B$ une nouvelle
relation que je n'aurois pû déterminer qu'en répétant l'experience \struck{et}
examinant avec soin la fig. \struck{la mesurant et
mesurant avec exactitude la distance détermi} tenant compte dela
\struck{distance} valeur de $x$ \struck{et} de $y$ \add{pour les} \struck{aux} points \add{de limites} auxquels \struck{comme} se terminent
\struck{aux par \ill{} limites de la plaque quarrée} les quatre lignes nodales
qui remplacent ici le\struck{s} \struck{\uncertain{qu}} systême des deux diagonales nodales.

Tout ce que l'on peut voir sans calcul, est que, \struck{si \uncertain{p}} lorsque l'on aura
ni $b=-B$, ni $b=0$ ou $B=0$, la figure nodale devra être
composée de quatre lignes \struck{nodales} placés à une distance
\note{En tête, au milieu, « 71 », sans parenthèses, barré d'un trait ; en
haut à droite, à l'encre, « 235 ». Le feuillet est collé bas sur sa
feuille de montage : le premier tiers de la vue est blanc. La première
ligne achève la phrase de la vue 478 (« et tandis que / $z'$
satisfera ») ; « a la condit. » y est barré. La ligne « car sans cette
condition, elle ne conviendroit qu'aux plaques dont les côtés seroient
appuyés. » est écrite en plus petit dans l'interligne, au-dessus de
« Cela posé », sans signe d'appel ; son dernier mot passe au bord du
feuillet et n'en montre que « app- ». « $\Pi''=1-9$ » est écrit tel quel ;
la lettre est le grand $\Pi$ des vues 476–478. Les capitales $A''$, $B$
sont droites. Les points de suspension des formules sont les siens. Les
numéros de figures sont ceux du traité de Chladni. Dans « du facteur »
(ligne des diagonales), la fin de « du » est surchargée. À la fig. 70 :
« et » est barré en fin de ligne ; « la mesurant et » et, à la ligne
suivante, « mesurant avec exactitude la distance détermi » sont barrés ;
« distance » est barré en tête de la ligne suivante, « et » est barré
entre $x$ et « de $y$ » ; « aux » est barré et « pour les » écrit
au-dessus, « de limites » au-dessus de « points auxquels » ; « comme »
est barré devant « se terminent ». La ligne « aux par … limites de la
plaque quarrée » est barrée entière ; un mot court y est surchargé et ne
se lit pas. « les » est corrigé en « le », suivi d'un « qu » barré, devant
« systême ». Devant « lorsque », « si » et une lettre, lue avec doute
« p », sont barrés. « nodales » est barré devant « placés ». La phrase se
poursuit à la vue suivante.}

\page{480}

de l'origine, dépendante dela relation particulière établie entre
$b$ et $B$. En effet, si on fait, par exemple, $\text{cos}\left(\frac{\pi y}{A}\right)=0$,
une ligne nodale commencera sur le côté parallelle à l'axe
des $\underline{y}$ lorsque \struck{l'équation} $\text{cos}^2\left(\frac{\pi x}{A}\right) = \frac{b+B}{4b}$ \add{sera satisfaite}; et comme
\struck{les deux angles qui diffèrent entr'eux de $\pi$ dont la somme}
on a une même valeur de $\underline{x}$ de part et d'autre de l'axe
des $\underline{y}$, et que les valeurs négatives de $\underline{x}$ prises \struck{de part et}
sont pareillement doubles, il est clair que l'on \struck{a quatre} \add{aura}
points qui satisfont àla même condition. On verra
également que la condition $\text{cos}^2\left(\frac{\pi y}{A}\right) = \frac{b+B}{4B}$ \struck{sera} \add{pourra être} remplie
par quatre autres points pris sur le côté parallelle à
l'axe des $\underline{x}$.

Lorsque $b=-B$, les points qui satisfont aux conditions
$\text{cos}^2\left(\frac{\pi x}{A}\right)=0$, et $\text{cos}^2\left(\frac{\pi y}{A}\right)=0$ se confondent, et il ne doit
y avoir, comme il n'y a en effet, que deux lignes nodales
diagonales. \add{Lorsqu'une des conditions \struck{sera} impossible, \struck{cequi doit arriver toutes les fois que} \add{c'est adire lorsque} $b$ et $B$ seront de signes differens sans que leur somme soit nulle; et aussi lorsque \add{l'un des coëfficiens sera $<$ que le tiers de l'autre} \struck{\uncertain{B par exemple sera $>$ que 3b}}, il ne devra y avoir que deux lignes nodales sinueuses dans \struck{une} direction parallelles à un des axes. \struck{Lorsque $b=0$ ou $B=0$ il n'existe qu'une seule condition} Lorsque l'un des coëfficiens $b$ ou $B$ sera nulle, ces lignes seront de simples droites. elles passeront par la suite des points qui rempliront l'une des deux conditions $\text{cos}\left(\frac{\pi.x}{A}\right)=\frac{1}{2}$, ou $\text{cos}\left(\frac{\pi.y}{A}\right)=\frac{1}{2}$; \struck{et il ne doit y avoir que deux lignes nodales parallelles} c'est adire qu'elles se trouveront distantes du côté \struck{le} plus proche \struck{de la plaque} auquel elles seront parallelles, d'un tiers de l'intervalle entre cette extremité et l'origine. \struck{ne \ill{} des axes.} Les lignes dont il s'agit \struck{serviront} rempliront les conditions des lignes fictives que j'ai emploiées dans l'application de l'intégrale (C), car, en les supposant tracés, elles serviront d'axes aux lignes sinueuses qui passeront par leurs intersections. desorte que ces intersections seront de véritables \struck{lignes} points invariables de repos.}

Dans tout autre cas+ il y aura nécessairement quatre lignes
nodales de direction plus ou moins inclinés à l'un
et l'autre \struck{des} axes.
\add{+ c'est adire lorsque les deux equations $\text{cos}^2\left(\frac{\pi x}{A}\right) = \frac{b+B}{4b}$, $\text{cos}^2\left(\frac{\pi y}{A}\right) = \frac{b+B}{4B}$ seront satisfaites.}
\note{En tête, au milieu, « (72) », non barré ; aucun nombre en haut à
droite : c'est le dos du feuillet de la vue 479. La première ligne
achève la phrase de la vue 479 (« placés à une distance / de
l'origine »). Dans la première formule, une lettre empâtée précède $y$
au numérateur, et le dénominateur $A$ est surchargé. « l'équation » est
barré ; « sera satisfaite » est écrit au-dessus de la fraction, d'une
plume qui la raye en partie. Dans « $\text{cos}^2(\pi x/A)$ » de la
même ligne, une lettre est noircie entre $\pi$ et $x$. La ligne « les
deux angles qui diffèrent entr'eux de $\pi$ dont la somme » est barrée
entière. « de part et » est barré en fin de ligne ; « a quatre » est
barré et « aura » écrit au-dessus, sans que « quatre » soit récrit.
« satisfont » est écrit sur une autre fin de mot ; « sera » est barré et
« pourra être » écrit au-dessus. Le long passage donné ici en addition,
de « Lorsqu'une des conditions » à « invariables de repos », est écrit
d'une plume plus fine et d'une écriture plus petite, d'abord dans
l'interligne après « diagonales. », puis sur les lignes laissées libres
au bas de la page et, pour sa fin (« qui passeront par leurs
intersections. desorte que ces intersections seront de véritables points
invariables de repos. »), dans la marge de gauche, face à « Dans tout
autre cas » ; il est donné ici d'un seul tenant, dans l'ordre où la
syntaxe le lit. « sera » est surchargé devant « impossible » ; « cequi
doit arriver toutes les fois que » est barré, et « c'est adire lorsque »
écrit au-dessus. Le groupe barré après « lorsque », sous « l'un des
coëfficiens sera < que le tiers de l'autre », est lu avec doute ; « une »
est barré devant « direction ». Dans « $\text{cos}(\pi.x/A) = \frac{1}{2}$ ; ou $\text{cos}(\pi.y/A) =
\frac{1}{2}$ », les fractions sont petites et mal formées ; la lecture $\frac{1}{2}$ est
offerte d'après la suite (« d'un tiers de l'intervalle »), non d'après
un calcul. Dans la ligne barrée après « l'origine », un mot entre « ne »
et « des axes » ne se lit pas. « serviront » est barré devant
« rempliront » ; « lignes » est barré devant « points ». Le grand D de
« Dans tout autre cas » est écrit sur le début de « qui passeront » de
l'addition. L'appel « + » après « cas » renvoie à la note « + c'est
adire lorsque … seront satisfaites. », écrite en petit sous la ligne.
Dans « y aura », la fin est surchargée. « des » est surchargé et barré
devant « axes ». Le texte s'arrête ici, aux deux tiers du feuillet, sur
un point : la fin d'un alinéa, peut-être d'une section ; la suite n'est
pas dans ce lot.}

\end{document}
